\chapter{Coffee Consumption Mathematics}

Coffee extraction occurs when hot water is poured over coffee grounds, causing
desirable compounds such as caffeine, carbohydrates, lipids, melanoidins and
acids to be extracted from the grounds.

Extraction yield refers to the solubles dissolved during brewing. This is often
expressed as a percentage of the coffee's mass. It is also known as solubles
yield or simply extraction. The extraction yield percentage describes the mass
transferred from coffee grounds to water, expressed as a percentage of the
initial mass of the grounds. It is given by eq.~\eqref{ch1:eq:extraction_yield}.

\begin{equation}
    Y = \frac{M_2 \cdot t}{M_1}
    \label{ch1:eq:extraction_yield}
\end{equation}

where $Y$ is the extraction yield expressed as a percentage, $t$ is the total
dissolved solids expressed as a percentage of the final beverage, $M_1$ is the
mass of the grounds in grams, and $M_2$ is the final beverage's mass in grams.
This means that an extraction yield of 20\% can be obtained by brewing 18 grams
of coffee, resulting a 36-gram final beverage with a $t$ of 10\%. Yield can also
be expressed as total dissolved solids, or parts-per-million (ppm).

Here is a subscript $x_1$ and superscript $x^2$. Here is a subsubscript $x_{y_1}$,
a supersuperscript $x^{y^1}$, and the others $x_{y^2}$ and $x^{y_2}$.